\section{五行棋盘的显式生成函数}\label{sec:five}
五行情形可以把一般矩阵公式压缩为六维多项式恒等式。
为使此实例可手工复核，本节给出证书和全部边界向量，
附录 \ref{app:certificate} 列出单列转移的构造。

用五位掩码表示行集合，最低位对应顶行。
偶数列可达状态的行反转轨道依次为
\begin{equation}\label{eq:even-orbits}
 \{0\},\{3,24\},\{6,12\},\{9,18\},\{15,30\},\{27\}.
\end{equation}
计数在行反转下相同。令 $u_j$ 的六个分量表示第 $2j$ 列后
相应轨道中\emph{每个单独状态}的计数，非轨道计数之和。
则 $u_0=\evec$，$u_{j+1}=Mu_j$，其中
\begin{equation}\label{eq:M5}
 M=\begin{pmatrix}
8&6&4&2&2&1\\
3&4&1&1&1&1\\
2&1&2&1&1&0\\
1&1&1&2&1&0\\
1&1&1&1&1&0\\
1&2&0&0&0&1
\end{pmatrix}.
\end{equation}
设
\begin{equation}\label{eq:PN}
\begin{split}
 P(t)&=1-15t+32t^2-15t^3+t^4,\\
 N(t)&=\bigl(1-7t+7t^2-t^3,\ 
 3t-4t^2+t^3,\ 2t-5t^2+t^3,\ t+t^2,\ t-t^3,\ t-t^3\bigr)^{\trans}.
\end{split}
\end{equation}
直接多项式乘法给出
\begin{equation}\label{eq:certificate}
 (I-tM)N(t)=P(t)\evec.
\end{equation}
由于 $P(0)=1$，该恒等式证明
\[
 f(t):=\sum_{j\ge0}u_jt^j=\frac{N(t)}{P(t)}.
\]
附录还把式 \eqref{eq:certificate} 写成四个系数向量等式，
因此不需要将符号软件输出作为黑箱。

左右反射在偶数、奇数半长度的终端行向量分别为
\begin{align*}
 w_{v,e}&=(1,2,2,2,2,1),&w_{v,o}&=(8,6,4,4,2,1),
\end{align*}
半周旋转相应为
\[
 w_{r,e}=(1,0,0,0,0,1),\qquad w_{r,o}=(2,2,0,0,0,1).
\]
偶数终端向量分别累计所有状态，或仅累计行反转固定状态；
奇数终端向量先作一次单列转移再累计。
引理 \ref{lem:half} 因而给出
\begin{equation}\label{eq:terminal5}
 V_5(x)=(w_{v,e}+xw_{v,o})f(x^2),\qquad
 R_5(x)=(w_{r,e}+xw_{r,o})f(x^2).
\end{equation}

\begin{theorem}[五行显式公式]\label{thm:five}
令 $D_5,H_5,V_5,R_5$ 分别为五行四种固定点数的生成函数，则
\begin{align}
 D_5(x)&=\frac{1-7x+7x^2-x^3}{P(x)},&
 H_5(x)&=\frac{1-x}{1-3x+x^2},\label{eq:five-DH}\\
 V_5(x)&=\frac{1+8x+8x^2-23x^3-9x^4+16x^5-x^7}{P(x^2)},\label{eq:five-V}\\
 R_5(x)&=\frac{1+2x-6x^2-7x^3+7x^4+6x^5-2x^6-x^7}{P(x^2)}.\label{eq:five-R}
\end{align}
轨道生成函数为
\begin{equation}\label{eq:five-A}
 A_5(x)=\frac{D_5(x)+H_5(x)+V_5(x)+R_5(x)}4.
\end{equation}
\end{theorem}
\begin{proof}
普通完整铺法对应空终端状态，故 $D_5=f_0$，
给出 \eqref{eq:five-DH} 的第一式。
将 \eqref{eq:PN} 代入 \eqref{eq:terminal5} 给出
\eqref{eq:five-V} 和 \eqref{eq:five-R}。
由引理 \ref{lem:h-odd}，$h_{5,n}=t_2(2n)$。
两行铺法从最左列分类，或者首列由一块竖直砖覆盖，
或者前两列由两块水平砖覆盖，所以
$t_2(j)=t_2(j-1)+t_2(j-2)$，初值为 $t_2(0)=t_2(1)=1$。
其两步转移矩阵为
$\left(\begin{smallmatrix}2&1\\1&1\end{smallmatrix}\right)$，
满足 $L^2-3L+I=0$。故偶长子序列 $b_n=t_2(2n)$ 的初值为 $1,2$，
且 $b_{n+2}=3b_{n+1}-b_n$，得到 $H_5$。
最后由定理 \ref{thm:odd} 或轨道平均公式得 \eqref{eq:five-A}。
\end{proof}

令 $Q(x)=(1-3x+x^2)P(x)P(x^2)=\sum_{j=0}^{14}q_jx^j$。
通分后的分子记为 $U(x)=\sum_{j=0}^{13}b_jx^j$。
表 \ref{tab:five-coeff} 给出二者全部系数和足以启动递推的 14 个初值。
因此
\begin{equation}\label{eq:rec5}
 a_5(n)=-\sum_{j=1}^{14}q_j a_5(n-j),\qquad n\ge14.
\end{equation}
本文只主张这一十四阶递推，不另行声称阶数的最小性。

\begin{table}[htbp]
\centering
\caption{五行分子、分母的升幂系数及递推初值；$b_{14}=0$。}\label{tab:five-coeff}
\begin{tabular}{rrrr}
\toprule
$j$ & $b_j$ & $q_j$ & $a_5(j)$\\
\midrule
0 & 1 & 1 & 1\\
1 & -13 & -18 & 5\\
2 & 6 & 63 & 33\\
3 & 194 & 144 & 329\\
4 & -365 & -1060 & 3818\\
5 & -371 & 1296 & 46878\\
6 & 1304 & 1312 & 584386\\
7 & -594 & -3492 & 7318152\\
8 & -703 & 1312 & 91752831\\
9 & 717 & 1296 & 1150774799\\
10 & -154 & -1060 & 14434490195\\
11 & -36 & 144 & 181060839597\\
12 & 15 & 63 & 2271175917429\\
13 & -1 & -18 & 28489049801109\\
14 & 0 & 1 & \text{由递推确定}\\
\bottomrule
\end{tabular}
\end{table}


Mathar \cite[表 12]{mathar2013} 给出的五行、列数 $2,4,6,8,10$ 的轨道数
为 $5,33,329,3818,46878$，与这里得到的 $a_5(1),\ldots,a_5(5)$ 一致。
该比较在公式推导之后进行，没有用这些数据拟合矩阵或生成函数。
